% !PRINTED 170
\weberpage{170}

Es ist aber
\[
F(\tfrac12,\tfrac12,1\pm\varepsilon,x) = \frac{\Pi(\pm\varepsilon)}{[\Pi(-\tfrac12)]^2} \sum_{\nu=0}^{\infty} \frac{\Pi(\nu-\tfrac12)^2}{\Pi(\nu\pm\varepsilon)\Pi(\nu)} x^\nu,
\]
und man kann unter Hinzufügung eines konstanten Faktors als
zweites partikulares Integral auch folgendes annehmen:
\[
\frac{1}{2\varepsilon}\cdot\sum_{\nu=0}^{\infty} \frac{\Pi(\nu-\tfrac12)^2}{\Pi(\nu)}\Bigl[\frac{-1}{\Pi(\nu+\varepsilon)} + \Bigl(\frac{x}{1-x}\Bigr)^{-\varepsilon}\frac{1}{\Pi(\nu-\varepsilon)}\Bigr]x^\nu.
\]
Entwickelt man hier nach Potenzen von $\varepsilon$ und setzt dann
$\varepsilon=0$, so folgt:
\[
\sum_{\nu=0}^{\infty} \frac{\Pi(\nu-\tfrac12)^2}{\Pi(\nu)^2} x^\nu \Bigl[\frac{\Pi'(\nu)}{\Pi(\nu)} - \frac12 \log\frac{x}{1-x}\Bigr],
\]
wofür man mit Benutzung der \gesperrt{Gaussschen} Formel
\[
\frac{\Pi(\nu-\tfrac12)\Pi(\nu)}{\Pi(2\nu)} = \sqrt{\pi}\,2^{-2\nu}
\]
mit Abwerfung eines konstanten Faktors auch setzen kann:
\weq{12}{\sum_{\nu=0}^{\infty} \frac{\Pi(2\nu)^2}{\Pi(\nu)^4}\Bigl(\frac{x}{16}\Bigr)^\nu \Bigl[\frac{\Pi'(\nu)}{\Pi(\nu)} - \frac12 \log\frac{x}{1-x}\Bigr].}
Es ist aber für ein ganzzahliges positives $\nu$:
\[
\frac{\Pi'(\nu)}{\Pi(\nu)} = \Pi'(0) + 1 + \frac12 + \frac13 + \cdots + \frac1\nu,
\]
und wenn wir also zur Abkürzung setzen:
% [sic] In the second line the source omits the plus sign between \cdots and 1/\nu.
\weq{13}{\begin{aligned}
F(x) &= \sum_{\nu=0}^{\infty} \frac{\Pi(2\nu)^2}{\Pi(\nu)^4}\Bigl(\frac{x}{16}\Bigr)^\nu, \\
G(x) &= 2\sum_{\nu=1}^{\infty} \frac{\Pi(2\nu)^2}{\Pi(\nu)^4}\Bigl(1+\frac12+\frac13+\cdots\frac1\nu\Bigr)\Bigl(\frac{x}{16}\Bigr)^\nu,
\end{aligned}}
so ergeben sich die beiden partikularen Lösungen der Differentialgleichung
 (8), wenn man für das zweite eine lineare Kombination
von (10) und (12) nimmt:
\weq{14}{y_1 = F(x), \qquad y_2 = \frac12 G(x) - \frac12 \log\frac{x}{16(1-x)}F(x).}
Um nun aber die partikularen Integrale $K$ und $K'$ durch
diese Funktionen darzustellen, müssen wir das Verhalten von $K,K'$
für $q=0$ untersuchen, wofür zugleich $\varkappa^2=0$ wird.

Es ist aber für $q=0$ nach \S\,25 und \S\,42:
\[
q^{-1}\varkappa^2 = 16, \qquad 2K = \pi,
\]
%CONT>>
